\documentclass{article}
\usepackage[utf8]{inputenc}
\usepackage[T1]{fontenc}
\usepackage[margin=2cm]{geometry}
\usepackage{hyperref}
\sloppy
\begin{document}
\section*{Historical source}
Encyclopaedia Britannica scan, source row 257. Parse 3.5 extracted the text. Super 3.5 VL corrected it from an enlarged crop of each block. It can still contain recognition errors.
Page headers, footers, and page numbers are removed.

\subsection*{Untitled continuation (E0012)}
\textit{Continues an entry from an earlier page.}

in his power to provide for him otherwise, Swift left his patron (1694) in discontent; having previously taken his master's degree at Oxford, by means of a testimonial from Dublin, in which the words of disgrace were omitted. He was resolved to enter into the church, where his first preferment was only L. 100 a year, being the prebend of Kilroot in Connor; which some time afterwards, upon Sir William Temple's earnestly inviting him back to his house at Moorpark, he refused in savour of a clergyman far advanced in years and burdened with a numerous family. For this man he solicited the prebend, to which he himself inducted him.

In 1699 Swift lost his patron Sir William Temple, who left him a legacy in money, with the property of his manuscripts; and, on his death-bed, obtained for him a promise from the king of the first prebend that should become vacant at Westminster or Canterbury. That this promise might not be forgotten, Swift dedicated to the king the voluminous works with which he was entrusted, and for a while attended the court; but soon found his solicitations hopeless. He was then invited by the earl of Berkeley to accompany him into Ireland, where, after suffering some cruel disappointments, he obtained the livings of Laracor and Rathbeggin in the diocese of Meath; and soon afterwards invited over the unfortunate Stella, a young woman of the name of Johnson, whose life he contrived to embitter, and whose days, though he certainly loved her, we may confidently affirm that he shortened by his caprice.

This lady is generally believed to have been the daughter of Sir William Temple's reward; but her niece, a Mrs Hearn, assured Mr Berkeley, the editor of a volume of letters entitled Literary Relics, that her father was a merchant, and the youngest brother of a good family in Nottinghamshire; that her mother was the intimate friend of lady Gifford, Sir William's heir; and that the herself was educated in the family with his niece, the late Mrs Temple of Moorpark by Farnham*. This story would be entitled to the fullest credit, had not Mrs Hearn affirmed, in the same letter, that before the death of Sir William Temple, Mrs Johnson's little fortune had been greatly injured by the South-Sea bubbles, which are known to have injured no person till the year 1720: (See Company, II. c.) When one part of a narrative is so palpably false, the remainder will always be received with hesitation. But whether Miss Johnson was the daughter of Temple's Reward or of the friend of lady Gifford, it is certain that Sir William left her L.1000; and that, accompanied by Mrs Dingley, whose whole fortune amounted to an annuity of L.27 for life, she went, in consequence of Swift's invitation, to Laracor. With these two ladies he passed his hours of relaxation, and to them he opened his bosom; but they never resided in the same house, and did he see either without a witness.

In 1701 Swift published A Discourse of the contests and dissensions in Athens and Rome. It was his first work, and indeed the only which he ever expressly acknowledged. According to his constant practice he had concealed his name; but after its appearance, paying a visit to some Irish bishop, he was asked by him if he had read that pamphlet, and what its reputation was in London. Upon his replying that he believed it was very well liked in London; "Very well liked!" said the bishop with some emotion. "Yes, Sir, it is one of the finest tracts that ever was written, and bishop Burnet is one of the best writers in the world." Swift, who always hated Burnet with something more than political rancour, immediately questioned his right to the work, when he was told by the bishop that he was "a young man;" and still persisting to doubt of the justice of Burnet's claim, on account of the dissimilarity of the style of the

pamphlet from that of his other works, he was told that he was “a very positive young man,” as no person in England but bishop Burnet was capable of writing it. Upon which Swift replied, with some indignation, I am to assure your lordship, however, that bishop Burnet did not write the pamphlet, for I wrote it myself. And thus was he forced in the heat of argument to avow what otherwise he would have for ever concealed.

Early in the ensuing firing king William died; and Swift, on his next visit to London found queen Anne upon the throne. It was generally thought, upon this event, that the Tory party would have had the ascendant; but, contrary to all expectation, the Whigs had managed matters so well as to get entirely into the queen's confidence, and to have the whole administration of affairs in their hands. Swift's friends were now in power; and the Whigs in general, knowing him to be the author of the Discourse on the Contests, \&c. which was written in defence of king William and his ministers against the violent proceedings of the house of commons, considered themselves as much obliged to him, and looked upon him as fast to their party. But Swift thought with the Whigs only in the state; for with respect to the church his principles were always those of a Tory. He therefore declined any intimate connection with the leaders of the party, who at that time professed what was called low church principles. But what above all shocked him, says Mr Sheridan, was their inviting Deists, Freethinkers, Atheists, Jews, and Infidels to be of their party, under pretence of moderation, and allowing a general liberty of conscience. As Swift was in his heart a man of true religion, he could not have borne, even in his private character, to have mixed with such a motley crew. But when we consider his principles in his political capacity, that he looked upon the church of England, as by law established, to be the main pillar of our newly erected constitution, he could not, consistently with the character of a good citizen, join with those who considered it more as an ornament than a support to the edifice; and could therefore look on with composure while it was undermining, or could even open the gate to a blind multitude, to try, like Sampson, their strength against it, and consider it only as sport. With such a party, neither his religious nor political principles would suffer him to join; and with regard to the Tories, as is usual in the violence of factions, they had run into opposite extremes, equally dangerous to the state. He was therefore during the earlier part of the queen's reign of no party, but employed himself in discharging the duties of his function, and in publishing from time to time such tracts as he thought might be useful. In the year 1704 he published the Tale of a Tub, which, considered merely as a work of genius, is unquestionably the greatest which he ever produced; but the levity with which religion was thought to be there treated, raised up enemies to him among all parties, and eventually precluded him from a bishopric. From that period till the year 1708, he seems to have employed himself in folios; but he then gave effectually to the public The Sentiments of a Church of England man, the ridicule of atrology under the name of Bickerstaff, the Argument against Abolishing Christianity, and the defence of the Sacramental Test.

Soon after began the busy and important part of Swift's life. He was employed (1710) by the primate of Ireland to solicit the queen for a remission of the first fruits and twentieth parts to the Irish clergy. This introduced him to Mr Harley, afterwards earl of Oxford, who, though a Whig himself, was at the head of the Tory ministry, and in great need of an auxiliary as able as Swift, by whose pen he and the other ministers might be supported in pamphlets,


\section*{Source record}
\begin{description}
\item[Scan ID] \texttt{britannica:309f765b1ecc83e0a7f7762e44b4a927a7555443:train:257}
\item[Image file] \texttt{page-0257-012cf5a82a33.jpg}
\item[Image SHA-256] \texttt{012cf5a82a33e24b9fbf8a9def7970590e8f8c97141a57a7ae5b3f1fcbaadffb}
\item[Dataset] NationalLibraryOfScotland/encyclopaedia\_britannica\_illustrated at revision \texttt{309f765b1ecc83e0a7f7762e44b4a927a7555443} (CC0-1.0)
\item[Extraction] model \texttt{nvidia/nemotron-3.5-parse}, raw reply \texttt{extract/raw/page-0257-012cf5a82a33.json}
\item[Raw reply SHA-256] \texttt{44ab71df936f8ba8e85527059491549e5a6561400f44e10d3dc5137e73b8e2bd}
\item[Correction] model \texttt{nvidia/Nemotron-3.5-Super-VL}, block crops, raw replies \texttt{correct/raw/crop\_nothink/page-0257-012cf5a82a33-b*.json}; super\_vl: 7
\item[Running header] none
\end{description}
\end{document}
